Пусть дан неориентированный граф $G$.

\begin{definition} Кликой называется любой его полный подграф. Независимым множеством называется любое множество вершин, между которыми нет рёбер.
\end{definition}

\begin{definition} Вершинным покрытием называется множество вершин, в котором лежит хотя бы один конец любого ребра.
\end{definition}


Сформулируем так же задачи распознавания:

\begin{itemize}
    \item $\clique = \{ (G,k)\ |$ в графе $G$ есть клика из $k$ вершин$\}$
    \item $\indset = \{ (G,k)\ |$ в графе $G$ есть независимое множество (антиклика) из $k$ вершин $\}$
    \item $\vertexcover = \{ (G,k)\ |$ в графе $G$ есть вершинное покрытие из $k$ вершин$\}$
\end{itemize}

\begin{statement} $\clique \leqslant_p \indset$ и $\indset \leqslant_p \clique$
\end{statement}

\begin{proof}
$(G,k) \in \indset \Leftrightarrow (\overline{G}, k) \in \clique$ 
\end{proof}

\begin{theorem}
Задача $\indset$ является \npc
\end{theorem}

\begin{proof}
Докажем, что $3\sat \leqslant_p \indset$. Построим граф следующим образом: каждому вхождению литерала сопоставим вершину графа. Таким образом, всего будет $3k$ вершин, где $k$ ~--- число дизъюнктов в 3-КНФ. Вершины, соответствующие литералам из одного дизъюнкта, соединим рёбрами и получим треугольник для каждого дизъюнкта. Также соединим все противоположные литералы, например, $p$ и $\neg p$. На этом построение графа заканчивается, размер независимого множества возьмём равным $k$.

$\Rightarrow$ Если у формулы есть выполняющий набор, то в каждом дизъюнкте есть хотя бы один истинный литерал. Возьмём из каждого треугольника по одной вершине, соответствующей истинному литералу из выполняющего набора. 
Соответствующие вершины образуют независимое множество, поскольку эти литералы из разных треугольников (дизъюнктов), а литералы вида $p$ и $\neg p$ не могут быть истинны одновременно. Этих вершин как раз $k$. 

$\Leftarrow$ Если, наоборот, есть независимое множество из $k$ вершин (больше не может быть, потому что тогда будет пара вершин из одного треугольника), то эти вершины обязаны соответствовать различным треугольникам (дизъюнктам). А поскольку среди этих вершин нет одновременно литералов вида $p$ и $\neg p$, то можно взять такие значения переменных, при которых все задействованные литералы истинны. Эти значения и будут выполняющим набором.
\end{proof}

\begin{statement}
    Задача $\vertexcover$ является \npc
\end{statement}

\begin{proof}
    Будем сводить $\indset$ к $\vertexcover$. 

Заметим, что независимое множество и вершинное покрытие являются дополнениями друг друга:

\begin{itemize}
    \item Если $2$ вершины из дополнения вершинного покрытия соединены ребром, то соответствующее ребро не покрыто в вершинном покрытии, но это противоречие $\Rightarrow$ дополнение вершинного покрытия это независимое множество.
    \item Если какое-то ребро не покрыто вершиной из дополнения независимого множества, то тогда оба конца ребра не лежат в дополнении независимого множества $\Rightarrow$ эти вершины лежат в независимом множестве, но это противоречие $\Rightarrow$ дополнение это вершинное покрытие.
\end{itemize}

Таким образом, $\indset \leqslant_p \vertexcover$ и $(G,k) \mapsto (G,n-k)$, то есть $(G,k) \in \indset \Leftrightarrow (G,n-k) \in \vertexcover$, где $n:=|G|$.

\end{proof}